\documentclass[11pt,a4paper]{article}
\usepackage[margin=2.3cm]{geometry}
\usepackage{amsmath,amssymb,booktabs,enumitem,hyperref,xcolor}
\hypersetup{colorlinks,linkcolor=teal!60!black,urlcolor=teal!60!black,citecolor=teal!60!black}
\setlist{itemsep=2pt}
\title{\textbf{Spora: method, validation and references}\\[2pt]\large NASA Space Apps Challenge 2026 --- Field Shift: Adapting Farms with NASA Data}
\date{November 2026}
\begin{document}
\maketitle

\section{Overview}
Spora compares three crop rotations for one field: the farmer's current rotation, a soil-building
option and a water-saving option. It uses daily NASA POWER weather, monthly greenness from NASA HLS
satellite images and slope from NASADEM. The method has five parts:
\begin{enumerate}
\item the weather change between two 20-year periods;
\item which months the soil is bare;
\item four numbers per rotation;
\item rules that build the two alternatives;
\item a ranking set by the farmer's priority, adjusted by the field-health rules (pests, diseases, nutrients, soil).
\end{enumerate}
Every threshold below is a stated default in the code and can be changed. None of them are fitted to data.

\section{Data}
\begin{tabular}{@{}p{3.6cm}p{11.4cm}@{}}
\toprule
NASA POWER & Daily $T_{\max}$, $T_{\min}$, precipitation, solar radiation, relative humidity and 2 m wind since 1981, about 0.5$^\circ$ grid \cite{power}.\\
NASA HLS v2.0 & Harmonized Landsat--Sentinel-2 surface reflectance, 30 m; NDVI \cite{tucker1979}; Fmask removes cloud, shadow, snow and water \cite{claverie2018}.\\
NASADEM & 1 arc-second (about 30 m) elevation; mean slope of the field \cite{nasadem}.\\
FAO-56 & Reference evapotranspiration and crop coefficients $K_c$ \cite{fao56}.\\
Crop table & Sowing and harvest months, $K_c$, base temperature, warmth needed to establish, frost sensitivity (Landwirtschaftskammer NRW, FAO).\\
\bottomrule
\end{tabular}

\section{Weather change}
Let $Y_{\text{now}}$ be the last 20 complete years and $Y_{\text{then}}$ the 20 years before. For an index
$X_y$ computed for each year $y$ (hot days $T_{\max}\ge 30\,^\circ$C, frost days $T_{\min}<0\,^\circ$C, heavy-rain days
$P\ge 20$ mm, dry spells of 10 days or more with $P<1$ mm in April--September, and yearly rain), Spora reports
\[
\bar X_{\text{then}}=\frac{1}{20}\sum_{y\in Y_{\text{then}}}X_y,\qquad
\bar X_{\text{now}}=\frac{1}{20}\sum_{y\in Y_{\text{now}}}X_y ,
\]
and labels the change as more common, less common or about the same. These indices follow the ETCCDI
climate-extremes indices (SU30, FD, R20mm, CDD) \cite{zhang2011,ipcc2021}.

\paragraph{Reference evapotranspiration.} Daily $ET_0$ (mm/day) uses FAO-56 Penman--Monteith \cite{fao56}:
\[
ET_0=\frac{0.408\,\Delta\,(R_n-G)+\gamma\,\dfrac{900}{T+273}\,u_2\,(e_s-e_a)}{\Delta+\gamma\,(1+0.34\,u_2)} ,
\]
where $R_n$ is net radiation, $G\approx0$ is the soil heat flux at daily scale, $T$ the mean temperature,
$u_2$ the wind speed at 2 m, $e_s-e_a$ the vapour-pressure deficit, $\Delta$ the slope of the saturation
vapour-pressure curve and $\gamma$ the psychrometric constant. When some variables are missing, Spora uses
Hargreaves--Samani \cite{hargreaves1985}:
\[
ET_0=0.0023\,R_a\,(T+17.8)\sqrt{T_{\max}-T_{\min}} .
\]

\section{Bare months}
A rotation repeats every $T$ months, indexed $m=1,\dots,T$. Spora defines
\[
u(m)=\begin{cases}1 & \text{soil bare in month } m,\\ 0 & \text{soil covered by a crop or cover crop.}\end{cases}
\]
\begin{itemize}
\item \textbf{From the farmer's input:} a crop covers the soil from its second month until harvest.
\item \textbf{From the satellite (demo field):} $\text{NDVI}=(\rho_{\text{NIR}}-\rho_{\text{red}})/(\rho_{\text{NIR}}+\rho_{\text{red}})$, using the median of clear pixels in the field. A month is bare when $\text{NDVI}<0.3$. The 0.3 limit is a common bare-soil threshold, not a calibrated value. Months without a clear image fall back to the farmer's input.
\end{itemize}

\section{Four numbers per rotation}
With $P(m)$ the mean monthly rain and $ET_0(m)$ the mean monthly reference evapotranspiration over
$Y_{\text{now}}$, and $K_c(m)$ the crop coefficient of whatever grows in month $m$:
\begin{align}
N_1 &= 12-\frac{12}{T}\sum_{m=1}^{T}u(m) && \text{months covered per year}\\
N_2 &= 100\,\frac{\sum_m u(m)\,P(m)}{\sum_m P(m)} && \text{\% of rain falling on bare soil}\\
N_3 &= \frac{12}{T}\sum_{m=1}^{T}\max\!\bigl(0,\;K_c(m)\,ET_0(m)-P(m)\bigr) && \text{extra water needed, mm/yr}\\
N_4 &= \frac{12}{T}\sum_{m=1}^{T}h(m) && \text{earning months per year}
\end{align}
where $h(m)=1$ if a crop that can be sold or grazed grows in month $m$.

\paragraph{Why these numbers.}
\begin{itemize}
\item $N_1$ and $N_2$ stand in for erosion risk. In the (R)USLE, soil loss is proportional to rainfall erosivity multiplied by a cover factor $C$, and $C$ is built from the share of erosive rain that falls in each crop stage \cite{wischmeier1978,panagos2015}. $N_2$ is a simplified version: it uses rain depth instead of erosivity and treats each month as either bare or covered. Slope from NASADEM is reported next to it, not multiplied in.
\item $N_3$ is the FAO-56 single-coefficient crop water demand minus rain, summed only over months with a deficit. Water stored in the soil and effective rain are ignored, so $N_3$ is a conservative upper guide.
\end{itemize}

\section{Building the alternatives}
\paragraph{Soil-building.} Spora finds the longest bare window $[s,e]$ of at least 2 months. A cover crop $c$
sown in calendar month $k$ inside the window is accepted only if all three tests pass:
\[
\underbrace{\sum_{i=0}^{1}\text{GDD}_c(k+i)\ \ge\ \text{GDD}^{\text{est}}_c}_{\text{enough warmth}},\qquad
\underbrace{\#\{y: T_{\min}<0 \text{ within 42 d of sowing}\}<4}_{\text{frost (if sensitive)}},\qquad
\underbrace{\sum_{i=0}^{1}\max(0,K_c\,ET_0-P)\le 30\text{ mm}}_{\text{rain to establish}}
\]
with $\text{GDD}_c(k)=\sum_{\text{days}}\max\bigl(0,\tfrac{T_{\max}+T_{\min}}{2}-T_{\text{base},c}\bigr)$ \cite{mcmaster1997}.
Legumes and longer covers are preferred. In rotations with a brassica (for example rapeseed), brassica and
clover covers are left out, because they host clubroot and \emph{Sclerotinia}.

\paragraph{Water-saving.} If $N_3\ge 50$ mm, Spora takes the main crop $j$ with the largest water gap
$G_j=\sum_{m\in j}\max(0,K_cET_0-P)$. It replaces that crop with a crop $r$ that fits the same season, ripens with
the local warmth and satisfies
\[
G_r\le 0.8\,G_j .
\]

\section{Ranking}
There is no weighted score. The farmer's priority picks the number that orders the options:
\[
\text{save water: } \min N_3;\qquad
\text{build soil: } \min N_2,\ \text{then}\ \max N_1;\qquad
\text{earn / feed animals: } \max N_4 .
\]

\paragraph{Warnings.} A warning (heat in a crop's sensitive months, frost after sowing, a dry spell of 14 days or more,
heavy rain on bare soil) is shown when it happens in at least $\lceil 0.2\times 20\rceil=4$ of the last 20 years. It is
reported as a frequency, for example ``about one year in four''.

\section{Do pests and diseases enter the formulation?}
\textbf{Yes, but not as a term inside $N_1$--$N_4$.} Spora has no data to put a number on disease pressure,
so the farmer's field-health answers act as \emph{constraints} on how the alternatives are built and as
\emph{rules} on the ranking. In both cases Spora can say in one sentence why something changed. The rules come
from the agronomy text supplied by the team's agriculturalist \cite{agrotext}.
The same advice is well supported in the literature on break crops and short rotations
\cite{kirkegaard2008,angus2015,bennett2012}.

Let $\mathcal{M}$ be the main crops in cycle order and $f(\cdot)$ the botanical family. Spora defines:
\[
R=\sum_{i}\mathbf{1}\bigl[f(c_i)=f(c_{i-1})=\text{Poaceae}\bigr] \quad\text{(cereal after cereal, cyclic)},
\]
\[
B=\mathbf{1}\bigl[\exists\ \text{cereal cover crop followed by a cereal main crop}\bigr]\quad\text{(green bridge)}.
\]
If the farmer reports past fungal, viral or pest problems ($H=1$), the following rules apply:
\begin{enumerate}
\item \textbf{Soil-building:} cereal cover crops are excluded when the rotation contains cereals, so $B=0$ by construction.
\item \textbf{Water-saving:} the replacement crop $r$ may not share a family with the crops before and after it.
\item \textbf{Ranking:} any option $o$ with $R_o>R_{\text{current}}$ or $B_o=1$ is moved below the others, whatever its $N$ values.
\end{enumerate}
Answers about soil and nutrients do not change the numbers. They add written warnings and advisor questions:
\begin{itemize}
\item heavy clay: root rots, soil-borne mosaic viruses, cereal ground beetle;
\item high nitrogen: powdery mildew, cereal leaf beetle;
\item viral problems: sow later in autumn, after the aphid flight;
\item fungal problems: certified seed and resistant varieties;
\item little nitrogen: prefer legumes;
\item sandy, acidic or chalky soil: general agronomy, labelled as such.
\end{itemize}

A natural next step is to add climate-driven disease-risk indicators from the same POWER data, such as mild
autumns that favour aphids, and wet springs or winters that favour mildew and root rots. Each would be counted
as ``years out of 20'', like the existing warnings.

\section{How the formulation is validated}
The method is new as a combination, but every part follows an established standard. This section separates what
has been checked from what has not.

\subsection*{Checked}
\begin{itemize}
\item \textbf{Standard methods:} $ET_0$ follows FAO-56 \cite{fao56}. The weather indices are ETCCDI definitions \cite{zhang2011}. GDD uses the standard averaging method \cite{mcmaster1997}. The bare-soil share follows the logic of the RUSLE cover factor \cite{wischmeier1978,panagos2015}.
\item \textbf{Input data:} NASA POWER temperature and radiation have been checked against ground stations and found suitable for agricultural modelling \cite{white2008,white2011}. HLS is a validated, harmonized product \cite{claverie2018}.
\item \textbf{Internal check:} Penman--Monteith and Hargreaves $ET_0$ were computed side by side for the demo site and agree in size.
\item \textbf{Satellite check (demo field near J\"ulich):} the bare and covered months from HLS NDVI for 2023--2025 match the rotation the farmer entered (sugar beet, winter wheat, rapeseed, winter wheat).
\item \textbf{Plausibility:} the demo results move in the expected direction. A winter rye cover crop raises $N_1$ from 7.8 to 9.0 months and lowers $N_2$ from 34\% to 24\%. Replacing sugar beet with sunflower lowers $N_3$ from 134 to 107 mm/yr.
\item \textbf{Software tests:} automated tests check the formulas, the field-health rules, the three languages and the refusal of fields drawn on water.
\end{itemize}

\subsection*{Not yet validated}
\begin{itemize}
\item There is no field trial or yield data showing that the suggested rotations perform better.
\item $N_2$ has not been compared with measured soil loss or with erosion maps such as ESDAC RUSLE.
\item POWER describes an area of about 50 km and has not been compared with a local DWD station for this field.
\item The thresholds (NDVI 0.3, 30 mm establishment gap, 20\% water gain, 4 in 20 years) are expert defaults and have not been tuned.
\item Agronomists have not yet reviewed the crop table and field-health rules outside NRW.
\end{itemize}
\textbf{Planned:}
\begin{itemize}
\item compare POWER with DWD station data;
\item compare $N_2$ with the European RUSLE layers \cite{panagos2015};
\item check NDVI-based cover on more NRW field blocks;
\item have regional advisors review the rules;
\item run a pilot with farmers, comparing Spora's advice with advisors' recommendations.
\end{itemize}

\section{Literature supporting the approach}
\begin{itemize}
\item \textbf{Cover crops and soil cover:} cover crops reduce erosion and nitrate loss and raise soil carbon \cite{tonitto2006,poeplau2015}. This supports using $N_1$ and $N_2$ as soil-building targets.
\item \textbf{Rotation and disease:} break crops (legumes, oilseeds) after cereals reduce soil-borne diseases and raise the following wheat yield \cite{kirkegaard2008,angus2015}. Short rotations lead to yield decline through pathogen build-up \cite{bennett2012}. This supports the cereal-after-cereal count $R$ and the green-bridge rule $B$.
\item \textbf{Water demand:} the FAO-56 crop coefficient method is the international standard for crop water needs \cite{fao56}.
\item \textbf{Changing extremes:} heat, drought and heavy rain are changing over Europe \cite{ipcc2021}, which is why Spora compares two 20-year periods instead of using a single long-term mean.
\end{itemize}

\begin{thebibliography}{99}\small
\bibitem{fao56} Allen, R.G., Pereira, L.S., Raes, D., Smith, M. (1998). \emph{Crop evapotranspiration: Guidelines for computing crop water requirements.} FAO Irrigation and Drainage Paper 56, FAO, Rome.
\bibitem{angus2015} Angus, J.F., Kirkegaard, J.A., Hunt, J.R., et al. (2015). Break crops and rotations for wheat. \emph{Crop \& Pasture Science} 66, 523--552.
\bibitem{bennett2012} Bennett, A.J., Bending, G.D., Chandler, D., Hilton, S., Mills, P. (2012). Meeting the demand for crop production: the challenge of yield decline in crops grown in short rotations. \emph{Biological Reviews} 87, 52--71.
\bibitem{claverie2018} Claverie, M., Ju, J., Masek, J.G., et al. (2018). The Harmonized Landsat and Sentinel-2 surface reflectance data set. \emph{Remote Sensing of Environment} 219, 145--161.
\bibitem{hargreaves1985} Hargreaves, G.H., Samani, Z.A. (1985). Reference crop evapotranspiration from temperature. \emph{Applied Engineering in Agriculture} 1(2), 96--99.
\bibitem{ipcc2021} Seneviratne, S.I., et al. (2021). Weather and climate extreme events in a changing climate. In: \emph{Climate Change 2021: The Physical Science Basis}, IPCC AR6 WGI, Ch.~11. Cambridge University Press.
\bibitem{kirkegaard2008} Kirkegaard, J., Christen, O., Krupinsky, J., Layzell, D. (2008). Break crop benefits in temperate wheat production. \emph{Field Crops Research} 107, 185--195.
\bibitem{mcmaster1997} McMaster, G.S., Wilhelm, W.W. (1997). Growing degree-days: one equation, two interpretations. \emph{Agricultural and Forest Meteorology} 87, 291--300.
\bibitem{nasadem} NASA JPL (2020). \emph{NASADEM Merged DEM Global 1 arc second V001.} NASA EOSDIS LP DAAC. doi:10.5067/MEaSUREs/NASADEM/NASADEM\_HGT.001.
\bibitem{panagos2015} Panagos, P., Borrelli, P., Meusburger, K., et al. (2015). Estimating the soil erosion cover-management factor at the European scale. \emph{Land Use Policy} 48, 38--50.
\bibitem{poeplau2015} Poeplau, C., Don, A. (2015). Carbon sequestration in agricultural soils via cultivation of cover crops -- A meta-analysis. \emph{Agriculture, Ecosystems \& Environment} 200, 33--41.
\bibitem{power} NASA Langley Research Center. \emph{POWER Project: Prediction Of Worldwide Energy Resources}, daily point API. \url{https://power.larc.nasa.gov/}
\bibitem{tonitto2006} Tonitto, C., David, M.B., Drinkwater, L.E. (2006). Replacing bare fallows with cover crops in fertilizer-intensive cropping systems: a meta-analysis of crop yield and N dynamics. \emph{Agriculture, Ecosystems \& Environment} 112, 58--72.
\bibitem{tucker1979} Tucker, C.J. (1979). Red and photographic infrared linear combinations for monitoring vegetation. \emph{Remote Sensing of Environment} 8, 127--150.
\bibitem{white2008} White, J.W., Hoogenboom, G., Stackhouse, P.W., Hoell, J.M. (2008). Evaluation of NASA satellite- and assimilation model-derived long-term daily temperature data over the continental US. \emph{Agricultural and Forest Meteorology} 148, 1574--1584.
\bibitem{white2011} White, J.W., Hoogenboom, G., Wilkens, P.W., Stackhouse, P.W., Hoell, J.M. (2011). Evaluation of satellite-based, modeled-derived daily solar radiation data for the continental United States. \emph{Agronomy Journal} 103, 1242--1251.
\bibitem{wischmeier1978} Wischmeier, W.H., Smith, D.D. (1978). \emph{Predicting rainfall erosion losses: a guide to conservation planning.} USDA Agriculture Handbook 537.
\bibitem{agrotext} \emph{Agronomy text on winter cereals, crop rotation, pests and diseases} (sections 2.4.1, 2.5), supplied by the team's agriculturalist.
\bibitem{zhang2011} Zhang, X., Alexander, L., Hegerl, G.C., et al. (2011). Indices for monitoring changes in extremes based on daily temperature and precipitation data. \emph{WIREs Climate Change} 2, 851--870.
\end{thebibliography}
\end{document}
